\begin{satz}{Existenz einer Banach--Mazur--Abbildung}
            {ExistenzEinerBanachMazurAbbildung}
Seien $(E, \norm\cdot_E)$, $(F, \norm\cdot_F)$
endlichdimensionale normierte $\K$--Vektorräume mit $\dim_\K E = \dim_\K F$.\\
Dann existiert eine Banach--Mazur--Abbildung zwischen $(E, \norm\cdot_E)$ und
$(F, \norm\cdot_F)$.
\end{satz}
\begin{proof}
Da die Vektorräume $E, F$ endlichdimensional sind, gibt es bijektive lineare
Abbildungen von $E$ nach $F$ und diese, sowie deren Umkehrabbildungen, sind
stetig.
Damit ist $\Inv(E, F) \neq \emptyset$.\\
Sei $d\coloneqq d_{BM}\left( (E, \norm\cdot_E), (F, \norm\cdot_F) \right)$.\\
Aufgrund der Infimumseigenschaft existiert zu jedem $n \in \N$ ein
$S_n \in \Inv(E, F)$ mit 
$d + \frac 1 n > \opnorm {S_n} E F \opnorm {S_n\inv}FE$.\\
Sei $\folge S n \N$ eine Auswahl solcher Elemente und
$T_n \coloneqq  \frac 1 {\opnorm {S_n} E F} S_n$ für alle $n \in \N$.\\
Dann gilt offenbar: 
$T_n \in \Inv(E, F)$ und $T_n\inv = \opnorm{S_n} F E \cdot S_n\inv$
für alle $n \in \N$.\\
Somit gilt für alle $n \in \N$:\[
\opnorm{T_n\inv}FE=\opnorm{S_n}EF\cdot\opnorm{S_n\inv}FE<d+\frac1n\le d+1
\text.\]
Nun ist auch $\folge T n \N \in \kran {L(E, F)} 0 1 ^\N$.\\
Da $E$ und $F$ endlichdimensional sind, ist auch $L(E, F)$ endlichdimensional
und somit ist $\kran{L(E,F)}01$ kompakt,
also existiert eine
Teilfolge $\left(T_{n_j}\right)_{j \in \N}$ von $\folge T n \N$, die gegen ein
$T_0 \in \kran{L(E, F)}01$ konvergiert.\\
Die Folge $(T_{n_j}\inv)_{j\in \N}$ ist offenbar ebenfalls
beschränkt.\\
Da $(E, \norm\cdot_E)$ als endlichdimensionaler normierter Raum vollständig
ist, erhalten wir mit~\ref{InvertierbarkeitDesGrenzwertsInvertierbarerOperatoren} die Konvergenz von
$\left(T_{n_j}\inv \right)_{j\in \N}$ und es gilt
\[T_0\text{ ist invertierbar mit }T_0\inv=\Limes j \infty T_{n_j}\inv\text. \]
Nach Konstruktion gilt für alle $n \in \N$\[ 
\abs{\opnorm {T_n} E F \opnorm{T_n\inv} F E - d} 
=\opnorm{T_n}EF\opnorm{T_n\inv}FE-d=\opnorm{T_n\inv}FE-d<\frac 1n\text.\]
Daher gilt\[
\opnorm {T_0\inv} F E = \Limes j \infty \opnorm {T_{n_j}\inv} F E = d\text.\]
Wegen $T_0 \in \kran{L(E, F)}01$ ist $T_0$ also eine Banach--Mazur--Abbildung.
\end{proof}